\footnotesize%
\pstart%
\noindent%
\pagelabel{aestimavirium_dat}%
(1)~Der vorliegende Textkomplex N.\;\ref{aestimavirium_intro} umfasst die umfangreiche Aufzeichnung N.\;\ref{41212} \textit{De aestimatione virium quas per aeris compressionem acquirere licet} und die weitgehend kürzere, titellose Aufzeichnung N.\;\ref{41213}, welcher editorisch der Titel \textit{Die Kraft der zusammengedruckten Luft} zugewiesen wurde.
Beide Texte sind der gleichen Aufgabe gewidmet, die elastische Kraft einer Luftmasse, die in einem verschlossenem Behälter zusammengedruckt wird, unter Berücksichtigung der Unterscheidung zwischen "toten" und "lebendigen" Kräften zu berechnen.
Hinsichtlich ihrer Textart unterscheiden sich N.\;\ref{41212} und N.\;\ref{41213} erheblich voneinander: \textit{De aestimatione virium quas per aeris compressionem acquirere licet} ist eine im Latenischen verfasste und auf zwei ganzen Foliobogen ausgedehnte ausführliche Untersuchung, die sich inhaltlich durch zahlreiche Wedungen und wiederholte Neuansätze auszeichnet, wohingegen \textit{Die Kraft der zusammengedruckten Luft} den Charakter eines gradlinigen und komprimierten Entwurfes aufweist, der unprätenziös auf Deutsch startet und nur ein kleines Quartoblatt in Anspruch nimmt.
Ihre inhaltliche Verwandschaft rechtfertigt jedoch die Entscheidung der Herausgeber, beide Texte einheitlich als Unterstücke zu edieren.
Denn N.\;\ref{41212} und N.\;\ref{41213} teilen nicht nur die gleiche oben genannte Zielsetzung, sondern auch die gleichen empirischen Voraussetzungen:
Untersucht wird in beiden Texten das Verhalten einer Luftmasse mit anfänglichem Volumen eines Kubikfußes, deren elastische Kraft unter gewöhnlichen Druckbedingungen als Fähigkeit geschätzt wird, ein Gewicht von 2.100 Pfund zu stemmen, (vgl. \refpassage{LH_35_09_25_001-002+004-005_001r_nld-1}{LH_35_09_25_001-002+004-005_001r_nld-2} mit \refpassage{LH_35_09_25_003_003r_nld-1}{LH_35_09_25_003_003r_nld-2}).
Vor allem aber in ihrem Ergebnis stimmen N.\;\ref{41212} und N.\;\ref{41213} überein:
Der besagten Luftmasse, in ein Zentel ihres ursprünglichen Volumens zusammengedrängt, komme eine elastische Kraft zu, die der kinetischen Kraft eines Körpers entspreche, das 4.200 Pfund wiege und aus einer Höhe von 0,9 Fuß falle %
(vgl. \refpassage{LH_35_09_25_001-002+004-005_005v_nld-1}{LH_35_09_25_001-002+004-005_005v_nld-2} mit \refpassage{LH_35_09_25_003_003v_nld-1}{LH_35_09_25_003_003v_nld-2} u. der Randbemerkung zum Diagramm \bracks{\textit{Fig.\:4}}, S.\;\pageref{LH_35_09_25_003_003v_Fig4}).
\pend%
\pstart%
Zu diesem gemeinsamen Ergebnis gelangt N.\;\ref{41212} erst am Ende eines langen und verwinkelten Weges, N.\;\ref{41213} hingegen ohne Umwendungen.
Aus diesem Unterschied lässt sich im Hinblick auf die relative Chronologie der zwei Texte schließen, dass N.\;\ref{41213} \textit{nicht} vor N.\;\ref{41212} entstanden sein kann.
Vielmehr dürfte N.\;\ref{41213} eine schnelle Zusammenfassunug der in N.\;\ref{41212} durchgeführten und abgeschlossenen Untersuchung sein, oder aber eine Art Zuarbeit, die ab einem bestimmten Punkt parallel zu N.\;\ref{41212} verfasst wurde und möglicherweise zur Orientierung bei der komplexen Untersuchung dienen sollte.
In beiden Fällen ist anzunehmen, dass N.\;\ref{41212} zumindest teilweise bereits vorlag, als N.\;\ref{41213} verfasst wurde; und auch in dem Fall, dass N.\;\ref{41213} auf N.\;\ref{41212} folgte, lag zwischen ihnen wahrscheinlich kein relevanter Zeitabstand.
Für eine gleichzeitige oder wenigstens zeitnahe Entstehung sprechen 
ferner die Überlieferungsumstände:
Im Leibniz-Nachlass findet sich N.\;\ref{41213}, wie schon die Blattzählung zeigt, in den gefalteten Bogen eingelegt, die als Textträger für N.\;\ref{41212} dienen (siehe \refpassage{LH_35_09_25_003_ubrlf-1}{LH_35_09_25_003_ubrlf-2}).
Diese Anordnung der Handschriften könnte auf Leibniz selbst zurückgehen.
\pend%
\pstart%
(2)~Ausschlaggebend für die Datierung der Aufzeichnung N.\;\ref{41212} \textit{De aestimatione virium quas per aeris compressionem acquirere licet} ist zunächst das in beiden Textträgern anzutreffende Wasserzeichen (vgl. LH~XXXV~9,\;25 Bl.\;\href{\Scan{LH_35_09_25_002wz}}{2} u.\;\href{\Scan{LH_35_09_25_005wz}}{5}): ein Posthorn an gekröntem Zierschield, das aus der Papiermühle in Hehlen\indexO{Hehlen (bei Goslar)} bei Goslar\indexO{Goslar} stammen könnte. %
Das Wasserzeichen kommt im Leibniz-Nachlass nach heutigem Wissensstand insgesamt sieben Mal vor und weist zwei verschiedene Varianten auf: die eine mit dünnhalsigem, die andere mit dickhalsigem Posthorn.
Belegt ist die Variante mit dünnhalsigem Posthorn u.\,a. in dem auf den Anfang 1685 datierbaren Zeugen~C des im Herbst 1685 angefertigten Briefes Leibnizens an seinen Dienstherrn, Herzog Ernst August\indexN{\textls{Braunschweig-Lüneburg}, Ernst August von, Herzog und Kurfürst von Hannover 1679\textendash1698} (LH~XXIII 2,\;1 Bl.\;1--4; vgl. \textit{LSB} I,\;4 N.\;173),\cite{01404} sowie im Träger des Briefes, den Johann Georg Lange\indexN{\textls{Lange}, Johann Georg (1687)} am 9.\;(19.) September 1686 an Leibniz gesendet hat (LBr~523 Bl.\;5; vgl. \textit{LSB} I,\;4 N.\;246).\cite{01410}
Die Variante mit dickhalsigem Posthorn findet sich u.\,a. im Konzept~A des Briefes Leibnizens vom April 1687 an de Bourgeauville,\indexN{\textls{Bourgeauville}, de (französischer Gesandter in Celle 1688--1689)} französischen Gesandten in Celle\indexO{Celle} (LBr~183 Bl.\;12--13; vgl. \textit{LSB} I,\;4 N.\;526).\cite{01405}
Die vier weiteren im Leibniz-Nachlass bekannten Vorkommnisse des Wasserzeichens sind bisher nicht mit Sicherheit datiert.
Die in den Textträgern der Aufzeichnung N.\;\ref{41212} vorkommende Variante ist die für die Zeitspanne von Anfang 1685 bis September 1686 belegte mit \textit{dünnhalsigem} Posthorn.
Folglich liegt die Annahme nahe, dass N.\;\ref{41212} im Zeitraum zwischen dem Winter 1684/1685 und dem Herbst 1686 verfasst wurde, als Leibniz sich oft im Harz\indexO{Harz} aufhielt (\textit{Chronik}, S.\;73--79).\cite{01236}
Da es aber nach heutigem Wissensstand nicht auszuschließen ist, dass die Variante des Wasserzeichens mit dünnhalsigem Posthorn auch früher oder später im Umlauf war, ist auch eine etwas frühere oder etwas spätere Datierung von N.\;\ref{41212} möglich.
Sowohl über den Terminus post quem wie auch über den Terminus ante quem der Entstehung von N.\;\ref{41212} lassen sich im Folgenden Beobachtungen anstellen, die Datierung des Wasserzeichens tatsächlich ergänzen.
\pend%
\pstart%
Gegen Ende der Aufzeichnung N.\;\ref{41212} (\refpassage{LH_35_09_25_001-002+004-005_005r_Bernoulli-1}{LH_35_09_25_001-002+004-005_005r_Bernoulli-2}) ist ein präziser Verweis auf die \textit{Dissertatio de gravitate aetheris}\cite{01403} des Baseler\indexO{Basel} Mathematikers Jacob Bernoulli\indexN{\textls{Bernoulli}, Jacob 1655\textendash1705} anzutreffen.
Diese 1683 in Amsterdam\indexO{Amsterdam} erschienene Abhandlung dürfte Leibniz durch die Lektüre der ihr gewidmeten anonymen Rezension im Märzheft 1683 der \textit{Acta eruditorum} (S.\;106\textendash112)\cite{01406} zur Kenntnis genommen haben.
Da er in N.\;\ref{41212} aber auf eine Ausführung in Bernoullis Traktat Bezug nimmt, auf die in der anonymen Rezension nicht eingegangen wird, kann nicht diese letztere die Quelle dafür gewesen sein.
Wann Leibniz die \textit{Dissertatio de gravitate aetheris}\cite{01403} tatsächlich gelesen hat, lässt sich nach heutigem Wissensstand nicht genau bestimmen.
Das Buch war auf der Frankfurter\indexO{Frankfurt am Main} Michaelis-Messe 1682 (\textit{Catalogus universalis}, Bl.\;\bracks{D4\,r\textsuperscript{o}})\cite{02596} angekündikgt und auf der Leipziger\indexO{Leipzig} Ostermesse 1683 (\textit{Catalogus universalis}, Bl.\;B3\,v\textsuperscript{o})\cite{02597} ausgestellt worden, weshalb es frühestens ab dem Frühjahr 1683 auf dem Buchmarkt erhältlich gewesen sein dürfte (in beiden Katalogen ist der Titel auch ungenau angegeben: wohl das Zeichen für eine späte Auslieferung).
Im Briefwechsel wird die Abhandlung zum ersten und bisher einzigen Mal erst nach Leibnizens Italienreise\indexO{Italien} erwähnt, nämlich in seinem Brief an Jacob Bernoulli\indexN{\textls{Bernoulli}, Jacob 1655\textendash1705} vom 24.\;September/4.\;Oktober 1690; vgl. \textit{LSB} III,\;4 N.\;279, S.\;573.24\,f.\cite{01260} %
Leibniz dürfte erst einige Monate nach dem Erscheinen der anonymen Rezension imstande gewesen sein, Bernoullis Traktat in die Hände zu bekommen, d.\,h. wohl frühestens in der zweiten Hälfte 1683 oder vielmehr im Laufe des Jahres 1684.
Dieser ungenaue Zeitpunkt, der sich insgesamt aber nur begrenzt vom frühesten Beleg des Wasserzeichens absezt, soll hier als Terminus post quem der Datierung von N.\;\ref{41212} betrachtet werden.
\pend%
\pstart%
Den Versuch, die elastische Kraft der zusammengedruckten Luft mathematisch zu bestimmen, stellt Leibniz auch in anderen Texten an, etwa in den von ihm eigenhändig datierten Entwürfen \textit{De aeris resistentia elastica} (Juli 1686; \textit{LSB} VIII,\;3 N.\;27\textsubscript{1}\cite{02561}) und \textit{De vi elastica ad rationes geometricas revocata tentamen} (27.\;August 1689; ebd. N.\;27\textsubscript{2}\cite{02562}).
Diese beiden Texte, die sowohl genetisch wie auch inhaltlich eng miteinander verbunden sind (siehe die einlaeitende Erläuterung in \textit{LSB} VIII,\;3, S.\;306), gehen von einem theoretischen Ansatz aus, welcher dem der Aufzeichnung N.\;\ref{41212} \textit{De aestimatione virium quas per aeris compressionem acquirere licet} stark ähnelt, obwohl sie auf konkrete empirische Werte verzichten.
Beide Texte \textit{LSB} VIII,\;3 N.\;27\textsubscript{1}\cite{02561} und N.\;27\textsubscript{2} sind aller Wahrschinlichkeit nach später entstanden als N.\;\ref{41212} in diesem Band.
Weder \textit{De aeris resistentia elastica} noch \textit{De vi elastica ad rationes geometricas revocata tentamen} weisen die Unsicherheiten und Schwankungen auf, die hingegen die Aufzeichnung N.\;\ref{41212} auszeichnen.
Die Entwürfe \textit{LSB} VIII,\;3 N.\;27\textsubscript{1}\cite{02561} und N.\;27\textsubscript{2}\cite{02562} setzen vielmehr -- stillschweigend -- den Ansatz voraus, der sich in N.\;\ref{41212} durchgesetzt hat, und operieren mit ihm ohne Zögerungen weiter.
Daher ist anzunehmen, dass die Aufzeichnung N.\;\ref{41212} in diesem Band bereits abgeschlossen war, als \textit{LSB} VIII,\;3 N.\;27\textsubscript{1}\cite{02561} in Juli 1686 angefertigt wurde.
Eine Entstehung von N.\;\ref{41212} nach diesem Datum würde sich als unerklärlich erweisen.
Dass in \textit{LSB} VIII,\;3 N.\;27\textsubscript{1}\cite{02561} keine Erwähnung der Aufzeichnung N.\;\ref{41212} erfolgt, könnte sogar als Zeichen dafür gedeutet werden,
\ledpb%
dass zwischen beiden Texten bereits ein gewisser zeitlicher Abstand vorlag.
Die erste Häfte 1686 ist demgemäß als spätester Entstehungszeitpunkt von N.\;\ref{41212} anzusehen, was wiederum mit der Datierung des Wasserzeichens gut übereinstimmt.
\pend%
\pstart%
Der Vollständigkeit halber ist festzuhalten, dass der Schlussteil des unter Vorbehalt auf die Zeitspanne von Ende Januar 1683 bis zum Juli 1686 datierbaren Entwurfes \textit{De vibrationibus aeris tensi} inhaltliche Verwandschaft mit der in der Aufzeichnung N.\;\ref{41212} ausgeführten Untersuchung aufweist (vgl. \textit{LSB} VIII,\;3 N.\;21,\cite{01440} S.\;284.20--285.10).
Welches Abhängigkeitsverhältnis zwischen beiden Texten besteht, lässt sich nicht leicht bestimmen.
Da diese Frage aber keinen relevanten Beitrag zur Datierung von N.\;\ref{41212} verspricht, wird hier auf sie nicht weiter eingegangen.
\pend%
\pstart%
(3)~Da die Aufzeichnung N.\;\ref{41213} \textit{Die Kraft der zusammengedruckten Luft} -- wie oben unter Punkt~(1) festgestellt wurde -- weder vor der Aufzeichnung N.\;\ref{41212} noch wesentlich später abgefasst worden sein dürfte, ist auch für sie den gleichen Entstehungszeitraum anzunehmen wie für N.\;\ref{41212}.
Die Spanne zwischen 1684 (oder Herbst 1683) und der ersten Hälfte 1686 wird demgemäß als Datierung für den gesamten Textkomplexes N.\;\ref{aestimavirium_intro} vorgeschlagen.
\pend%
\vspace{4.0em}%
\selectlanguage{latin}%
\frenchspacing%
\normalsize%
